%% Reproducibility and Responsibility Checklist --- Paper 1: Where Does the Model Catch Its Own Mistakes?
%% Required for CAISc 2026 submission.

\section*{Reproducibility and Responsibility Checklist}

\subsection*{Reproducibility}

\begin{enumerate}

\item \textbf{Model and Code.}
\textit{Are the model weights and experimental code publicly available or will they be released?} \\
\textbf{Yes.} The base model (\texttt{deepseek-ai/DeepSeek-R1-Distill-Qwen-7B}) is publicly available via Hugging Face under the DeepSeek license. All experimental code---self-correction detection, activation extraction, probe training, head attribution, causal ablation harness, and figure generation---will be released as a public repository upon acceptance.

\item \textbf{Dataset Availability.}
\textit{Are the datasets used or constructed available?} \\
\textbf{Yes (derived dataset).} The self-correction dataset was derived from model generations on the publicly available GSM8K dataset~\citep{cobbe2021gsm8k} using the programmatic regex pipeline described in Section~\ref{sec:data} and Appendix~\ref{app:data}. GSM8K is available at \url{https://github.com/openai/grade-school-math} under the MIT license. The derived activation dataset (positional windows around correction markers) will be released alongside the code.

\item \textbf{Computing Requirements.}
\textit{What computational resources are required to reproduce the experiments?} \\
A single GPU with 48\,GB VRAM (e.g., NVIDIA RTX 6000 Ada) is required to load DeepSeek-R1-Distill-Qwen-7B at float16 precision via TransformerLens with activation storage. The activation extraction and probe training pipeline runs in approximately 2--4 hours. The causal ablation experiment (50 problems $\times$ 5 conditions) requires additional generation time of approximately 1--2 hours. Specific software requirements: CUDA 12.8+, Python 3.11, PyTorch 2.11, TransformerLens 3.2.1, Transformers 5.9.0, scikit-learn 1.4+. A \texttt{requirements.txt} with pinned versions will be provided.

\item \textbf{Hyperparameters and Random Seeds.}
\textit{Are all hyperparameters and random seeds fully specified?} \\
\textbf{Yes.} All random seeds for data splits, probe training, and head selection are fixed and reported (ablation generation uses stochastic sampling, so reported self-correction rates are point estimates):
\begin{itemize}
  \item Probe training: \texttt{random\_state=42} throughout.
  \item Random head selection (controls): a single \texttt{numpy.random.default\_rng(42)} stream, drawn sequentially --- random-3 = heads \{18, 2, 20\}, random-5 = heads \{20, 2, 27, 5, 18\}.
  \item Probe hyperparameters: \texttt{LogisticRegression(C=0.1, max\_iter=1000, random\_state=42)} on $z$-scored features.
  \item Train/validation split: 80/20 stratified, \texttt{random\_state=42}.
  \item Ablation generation: stochastic sampling (\texttt{temperature=0.7}, \texttt{top\_p=0.95}, \texttt{do\_sample=True}).
\end{itemize}

\item \textbf{Statistical Testing.}
\textit{Are statistical tests and confidence intervals reported?} \\
\textbf{Yes.} Fisher exact tests (two-sided) are used for all ablation comparisons. Reported values:
\begin{itemize}
  \item Top-3 vs.\ random-3: $p = 0.006$.
  \item Top-3 vs.\ baseline: $p = 0.39$ (not significant).
\end{itemize}
Effect sizes are reported as percentage-point differences in SC rates.

\item \textbf{Reproducibility Scope.}
\textit{Which claims are empirically supported, and which involve interpretation?} \\
All quantitative claims in Sections~\ref{sec:results} are derived directly from the experimental pipeline. The interpretations regarding the layer-0 caveat (AUROC 1.0 reflecting token-content differences), the distributed-circuit hypothesis (top-3 = top-5), and the mechanism behind random-$k$ ablation increasing SC rates are labeled as interpretations and marked as requiring further investigation.

\end{enumerate}

\subsection*{Responsibility}

\begin{enumerate}
\setcounter{enumi}{6}

\item \textbf{Dual-Use Risks.}
\textit{Could the methods or findings be misused?} \\
\textbf{Low risk.} The self-correction circuit identified here is a research characterization. Applying it as a steering or ablation tool requires white-box model access (weights, residual stream hooks), which is not available to typical API users. Suppressing self-correction in a reasoning model would degrade its accuracy, which is not a useful misuse vector. We see no immediate safety concerns beyond those that apply to mechanistic interpretability research generally.

\item \textbf{Limitations Acknowledged.}
\textit{Are the paper's limitations discussed openly?} \\
\textbf{Yes.} Section~\ref{sec:discussion} acknowledges: (a) the layer-0 probe peak reflects surface token content, not transformer computation; (b) the 4pp drop vs.\ baseline is not statistically significant in isolation; (c) the random-$k$ 100\% rate may reflect a confound rather than a true baseline; (d) the effect may be distributed across multiple heads; (e) results are from a single model.

\item \textbf{Broader Impacts.}
\textit{Are broader societal impacts considered?} \\
Understanding where and how reasoning models detect their own errors is foundational for model reliability research. If self-correction can be amplified via targeted steering, this could improve accuracy on difficult reasoning tasks without additional training. There are no anticipated negative societal impacts from this mechanistic research.

\item \textbf{Human Subjects / Privacy.}
\textit{Does the study involve human participants or personal data?} \\
\textbf{No.} All data consists of model-generated text and internal activations on publicly available math word problems (GSM8K). No human subjects are involved. No personal or private data is used or generated.

\end{enumerate}
